\documentclass[jp]{../../elegantbook}
% ========== 表紙と基本情報 ==========
\title{ElegantBook テスト文書}
\subtitle{テンプレートの核心機能}
\author{テストエンジニア}
\institute{テンプレート試験グループ}
\date{\today}
\version{v1.0}
\extrainfo{本ファイルはテンプレート機能試験専用であり、学術的価値を一切有しません}
\bioinfo{試験ロット}{TEST-FULL-20260506}
\logo{../../figure/logo-blue.png}
\cover{../../figure/cover.jpg}

% ========== 目次と基本設定 ==========
\setcounter{tocdepth}{3} % 目次を第3階層まで表示

% ========== 試験用参考文献 ==========
\addbibresource[location=local]{../en/reference.bib}

\begin{document}

\maketitle
\frontmatter
\tableofcontents

% ========== 序文（セクション分割付き） ==========
\pagenumbering{Roman}
\setcounter{page}{0}
\chapter*{試験序文}
\markboth{試験序文}{}
\section*{文書説明}
\markright{文書説明}
本文書は ElegantBook 全機能の試験ケースであり、テンプレートの基本機能・高級機能を全て網羅しています。本文中の内容はすべて架空の試験データであり、実質的な意味は持ちません。
\newpage
\section*{試験範囲}
\markright{試験範囲}
本試験の対象：表紙、目次、フォント、多階層リスト、複雑な数式、全ての数式環境、相互参照、図表、傍注、章概要、演習問題、参考文献、付録。

\mainmatter

% ========== 第1章：基本組版と多階層リスト ==========
\chapter{基本組版とリスト試験}
\begin{introduction}[本章試験項目]
\item 日本語フォントと文字書式
\item 3階層無順序リストの入れ子
\item 3階層順序リストの入れ子
\item 通常段落と強調文字列
\end{introduction}

\section{フォントと文字試験}
通常本文の試験。\textbf{太字}、\textit{斜体}、\underline{下線付き}、\textcolor{red}{赤色文字試験}。

\section{3階層無順序リスト試験}
\begin{itemize}
\item 第1階層試験項目 A
\item 第1階層試験項目 B
    \begin{itemize}
    \item 第2階層試験項目 B1
    \item 第2階層試験項目 B2
        \begin{itemize}
        \item 第3階層試験項目 B2-1
        \item 第3階層試験項目 B2-2
            \begin{itemize}
            \item 第4階層試験項目 B2-1a
            \item 第4階層試験項目 B2-2b
            \end{itemize}
        \end{itemize}
    \end{itemize}
\end{itemize}

\section{3階層順序リスト試験}
\begin{enumerate}
\item 第1階層順序試験 1
\item 第1階層順序試験 2
    \begin{enumerate}
    \item 第2階層順序試験 2-1
    \item 第2階層順序試験 2-2
        \begin{enumerate}
        \item 第3階層順序試験 2-2-1
        \item 第3階層順序試験 2-2-2
            \begin{enumerate}
            \item 第4階層順序試験 2-2-2-1
            \item 第4階層順序試験 2-2-2-2
            \end{enumerate}
        \end{enumerate}
    \end{enumerate}
\end{enumerate}

% ========== 第2章：複雑数式試験 ==========
\chapter{複雑な数式の試験}
\begin{introduction}
\item 単一行複雑数式
\item 複数行数式（align環境）
\item 行列・行列式
\item 多重積分・総和・極限・分数式
\item 数式ラベルと相互参照
\end{introduction}

\section{単一行複雑数式}
インライン複雑数式：$\sum_{i=1}^n \frac{x_i^2}{\sqrt{y_i+1}} \geq \bar{x}^2$、$\lim_{n\to\infty} \left(1+\frac{1}{n}\right)^n = e$。

\section{複数行数式と行列}
\begin{align}
a^2 + b^2 &= c^2 \label{eq:gougu} \\
\int_{0}^{1}\int_{0}^{1} f(x,y) dxdy &= 1 \label{eq:double} \\
\begin{vmatrix}
a & b \\
c & d
\end{vmatrix} &= ad-bc \label{eq:det}
\end{align}

\section{大型演算子数式}
\begin{equation}\label{eq:series}
\sum_{k=0}^{\infty} \frac{x^k}{k!} = e^x,\quad \oint_{C} Pdx+Qdy = \iint_{D}\left(\frac{\partial Q}{\partial x}-\frac{\partial P}{\partial y}\right)dxdy
\end{equation}

数式相互参照：三平方の定理\eqref{eq:gougu}、二重積分\eqref{eq:double}、行列式\eqref{eq:det}、級数\eqref{eq:series}。

% ========== 第3章：全数式環境試験 ==========
\chapter{全数式環境の試験}
\begin{introduction}
\item 定義・定理・補題・系・公理・公準
\item 命題・例題・問題・演習
\item 注意・備考・結論・仮定・性質
\item 解答・証明環境
\end{introduction}

\section{環境試験}

% 1. 定義
\ifdefined\simpleTest
\begin{definition}[試験定義]\label{def:test}
\else
\begin{definition}[試験定義]{test}
\fi
本定義環境は試験用内容であり、実際の数学的意味はなく、書式スタイルの検証のみを目的とします。
\end{definition}

% 2. 定理
\ifdefined\simpleTest
\begin{theorem}[試験定理]\label{thm:test}
\else
\begin{theorem}[試験定理]{test}
\fi
$x>0$ ならば $x+1>1$ が成り立ち、数式\eqref{eq:gougu}に対応します。
\end{theorem}

% 3. 補題
\ifdefined\simpleTest
\begin{lemma}[試験補題]\label{lem:test}
\else
\begin{lemma}[試験補題]{test}
\fi
任意の実数 $x$ に対し $x^2\geq0$ が成立します。
\end{lemma}

% 4. 系
\ifdefined\simpleTest
\begin{corollary}[試験系]\label{cor:test}
\else
\begin{corollary}[試験系]{test}
\fi
補題\ref{lem:test}より $(x-1)^2\geq0$ が導かれます。
\end{corollary}

% 5. 公理
\ifdefined\simpleTest
\begin{axiom}[試験公理]\label{axi:test}
\else
\begin{axiom}[試験公理]{test}
\fi
試験用公理であり、実質的な意味は持ちません。
\end{axiom}

% 6. 公準
\ifdefined\simpleTest
\begin{postulate}[試験公準]\label{pos:test}
\else
\begin{postulate}[試験公準]{test}
\fi
試験用公準であり、環境動作確認のみを行います。
\end{postulate}

% 7. 命題
\ifdefined\simpleTest
\begin{proposition}[試験命題]\label{pro:test}
\else
\begin{proposition}[試験命題]{test}
\fi
試験用命題内容で、定義\ref{def:test}と関連付けられています。
\end{proposition}

% 8. 例題
\begin{example}\label{exa:test}
例題環境の試験、自動番号付き。
\end{example}

% 9. 問題
\begin{problem}\label{probl:test}
問題環境の試験、解決すべき課題を記述する用途。
\end{problem}

% 10. 演習
\begin{exercise}\label{exer:test}
計算せよ：$1+2+3+\dots+10$。
\end{exercise}

% 11. 注意
\begin{note}
注意環境：番号なし、重要な補足提醒に使用。
\end{note}

% 12. 備考
\begin{remark}
備考環境：補足説明を記載する用途。
\end{remark}

% 13. 結論
\begin{conclusion}
結論環境：試験のまとめ内容を記載。
\end{conclusion}

% 14. 仮定
\begin{assumption}
仮定環境：試験条件の設定に使用。
\end{assumption}

% 15. 性質
\begin{property}
性質環境：対象の特性を記述。
\end{property}

% 16. 解答
\begin{solution}
演習\ref{exer:test}の解答：総和は $55$（試験回答）。
\end{solution}

% 17. 証明
\begin{proof}
定理\ref{thm:test}の証明：$x>0$ のため両辺に1を加えれば証明完了。
\end{proof}

% ========== 第4章：相互参照・図表・傍注試験 ==========
\chapter{相互参照と拡張機能}
\begin{introduction}
\item 全種類の相互参照
\item 簡単な図表試験
\item 傍注機能試験
\item 章ジャンプ動作検証
\end{introduction}

\section{全相互参照一覧}
定義\ref{def:test}、定理\ref{thm:test}、補題\ref{lem:test}、系\ref{cor:test}、公理\ref{axi:test}、公準\ref{pos:test}、命題\ref{pro:test}、例題\ref{exa:test}、問題\ref{probl:test}、演習\ref{exer:test}。

\section{図表試験}
\begin{figure}[h]
    \centering
    \includegraphics[width=0.5\textwidth]{../../image/scatter.jpg}
    \caption{試験図}\label{fig:test}
\end{figure}
\begin{table}[h]
    \centering
    \begin{tabular}{c*{6}{c}}
    \toprule
        $n$ & 1 & 2 & 3 & 4 & 5 & 6 \\
    \midrule
        $a_n$ & 1 & 1 & 2 & 3 & 5 & 8 \\ 
    \bottomrule
    \end{tabular}
    \caption{試験表}\label{tab:test}
\end{table}
図表参照：\figref{fig:test} は試験図、\tabref{tab:test} は試験表です。

% ========== 章末演習 ==========
\begin{problemset}[本章試験演習]
\item 定義\ref{def:test}の表示スタイルを検証せよ
\item 数式\eqref{eq:double}を導出せよ
\item 図\ref{fig:test}が正常に表示されるか確認せよ
\end{problemset}

% ========== 参考文献 ==========
\nocite{*}
\printbibliography

% ========== 付録 ==========
\appendix
\chapter{試験付録}
付録内容は付録章・ヘッダー・ページ番号・組版スタイルの検証用であり、実質的な情報は含みません。
\section{付録試験1}
付録内容は付録章・ヘッダー・ページ番号・組版スタイルの検証用であり、実質的な情報は含みません。
\newpage
\section{付録試験2}
付録内容は付録章・ヘッダー・ページ番号・組版スタイルの検証用であり、実質的な情報は含みません。
\end{document}